\subsection{Radicals, Hyperbolic Spaces, and Witt Theory}

Throughout this subsection, metric vector spaces are finite-dimensional over a fixed field $\mathbb{F}$.

\begin{definition} \label{def:metric-orthogonal-direct-sum} \hlt{Orthogonal Direct Sum}\\
For $\mathcal{S},\mathcal{T}\leq\mathcal{V}$, write $\mathcal{V}=\mathcal{S}\odot\mathcal{T}$ when $\mathcal{V}=\mathcal{S}\oplus\mathcal{T}$ and $\mathcal{S}\perp\mathcal{T}$; this is the \hlt{orthogonal direct sum} of $\mathcal{S}$ and $\mathcal{T}$.
\end{definition}

\begin{theorem} \label{thm:metric-orthogonal-complement-dimension} \hlt{Dimensions of Orthogonal Complements}\\
Let $\mathcal{S}\leq\mathcal{V}$.
\begin{enumerate}[label=(\roman*)]
\setlength{\itemsep}{0pt}
\item If $\mathcal{V}$ or $\mathcal{S}$ is nonsingular, then $\dim(\mathcal{S})+\dim(\mathcal{S}^{\perp})=\dim(\mathcal{V})$.
\item If $\mathcal{S}$ is nonsingular, then $\mathcal{V}=\mathcal{S}\odot\mathcal{S}^{\perp}$.
\item If $\mathcal{V}$ is nonsingular, then $\mathcal{S}^{\perp\perp}=\mathcal{S}$.
\item If $\mathcal{V}$ is nonsingular, then $\mathcal{S}$ is nonsingular $\Longleftrightarrow$ $\mathcal{S}^{\perp}$ is nonsingular.
\end{enumerate}
\begin{proof}
\begin{enumerate}[label=(\roman*)]
\setlength{\itemsep}{0pt}
\item[$(i)$] Define $\rho:\mathcal{V}\to\mathcal{S}'$ by $\rho(\mathbf{x})(\mathbf{s})=\innerprod{\mathbf{s},\mathbf{x}}$. Then $\rho\in\LL(\mathcal{V},\mathcal{S}')$ and $\ker(\rho)=\mathcal{S}^{\perp}$ by bilinearity and the definition of $\mathcal{S}^{\perp}$.\\
Fix $f\in\mathcal{S}'$. If $\mathcal{V}$ is nonsingular, choose $\mathcal{V}=\mathcal{S}\oplus\mathcal{C}$ by Proposition \ref{prop:complements-exist-finite-dimensions} and define $\widetilde f(\mathbf{s}+\mathbf{c})=f(\mathbf{s})$.\\
Then $\widetilde f\in\mathcal{V}'$ by the direct-sum decomposition. Fix $\mathbf{x}\in\mathcal{V}$ with $(\forall\,\mathbf{s}\in\mathcal{S})\,f(\mathbf{s})=\widetilde f(\mathbf{s})=\innerprod{\mathbf{s},\mathbf{x}}$ by Theorem \ref{thm:riesz-map-for-nonsingular-metric-spaces}.\\
If $\mathcal{S}$ is nonsingular, fix $\mathbf{x}\in\mathcal{S}$ with $(\forall\,\mathbf{s}\in\mathcal{S})\,f(\mathbf{s})=\innerprod{\mathbf{s},\mathbf{x}}$ by Theorem \ref{thm:riesz-map-for-nonsingular-metric-spaces} for $\mathcal{S}$. Thus $\rho$ is surjective in either case.\\
The identity $\dim(\mathcal{V})=\dim(\ker\rho)+\dim(\mathcal{S}')=\dim(\mathcal{S}^{\perp})+\dim(\mathcal{S})$ holds by Theorem \ref{thm:rank-nullity} and Theorem \ref{thm:dual-basis-and-dimension}(ii).
\item[$(ii)$] If $\mathcal{S}$ is nonsingular, then $\mathcal{S}\cap\mathcal{S}^{\perp}=\operatorname{rad}(\mathcal{S})=\{\mathbf{0}\}$. Thus $\mathcal{V}=\mathcal{S}\oplus\mathcal{S}^{\perp}$ by (i), and $\mathcal{V}=\mathcal{S}\odot\mathcal{S}^{\perp}$ by $\mathcal{S}\perp\mathcal{S}^{\perp}$.
\item[$(iii)$] Suppose that $\mathcal{V}$ is nonsingular. Then $\mathcal{S}\subseteq\mathcal{S}^{\perp\perp}$ by orthogonality.\\
The equality $\dim(\mathcal{S}^{\perp\perp})=\dim(\mathcal{V})-\dim(\mathcal{S}^{\perp})=\dim(\mathcal{S})$ holds by (i). Hence $\mathcal{S}^{\perp\perp}=\mathcal{S}$ by Proposition \ref{prop:dimension-bounds-and-basis-criteria}(ii).
\item[$(iv)$] ($\Rightarrow$) If $\mathcal{S}$ is nonsingular, then $\operatorname{rad}(\mathcal{S}^{\perp})=\mathcal{S}^{\perp}\cap\mathcal{S}^{\perp\perp}=\mathcal{S}^{\perp}\cap\mathcal{S}=\{\mathbf{0}\}$ by (iii).
\item[$(iv)$] ($\Leftarrow$) If $\mathcal{S}^{\perp}$ is nonsingular, then $\operatorname{rad}(\mathcal{S})=\mathcal{S}\cap\mathcal{S}^{\perp}=\mathcal{S}^{\perp\perp}\cap\mathcal{S}^{\perp}=\{\mathbf{0}\}$ by (iii).
\end{enumerate}
\hlt{Structural Commentary}: The restriction map has kernel $\mathcal{S}^{\perp}$; rank--nullity and double complementation settle (i)--(iv).\\
\hlt{Key Insight}: Nonsingularity makes orthogonal complementation dimension-reversing and involutive.
\end{proof}
\end{theorem}

\begin{theorem} \label{thm:radical-splitting-metric-space} \hlt{Splitting off the Radical}\\
Let $R=\operatorname{rad}(\mathcal{V})$. Then $\exists\,\mathcal{S}\leq\mathcal{V}$ such that $\mathcal{V}=R\odot\mathcal{S}$, $\operatorname{rad}(\mathcal{S})=\{\mathbf{0}\}$, and $R$ is totally singular.
\begin{proof}
Choose $\mathcal{S}\leq\mathcal{V}$ with $\mathcal{V}=R\oplus\mathcal{S}$ by Proposition \ref{prop:complements-exist-finite-dimensions}. Since $R\perp\mathcal{V}$, $\mathcal{V}=R\odot\mathcal{S}$ and the form on $R$ vanishes, so $\operatorname{rad}(R)=R$.\\
If $\mathbf{s}\in\operatorname{rad}(\mathcal{S})$, then $\mathbf{s}\perp\mathcal{S}$ and $\mathbf{s}\perp R$, whence $\mathbf{s}\in\operatorname{rad}(\mathcal{V})\cap\mathcal{S}=R\cap\mathcal{S}=\{\mathbf{0}\}$. Thus $\mathcal{S}$ is nonsingular.\\
\hlt{Structural Commentary}: A linear complement is automatically orthogonal to the radical and has zero restricted radical.\\
\hlt{Key Insight}: All degeneracy lies in a zero-form orthogonal summand.
\end{proof}
\end{theorem}

\begin{proposition} \label{prop:metric-quotients-by-radical-subspaces} \hlt{Metric Quotients}\\
Let $\mathcal{U}\leq\mathcal{V}$ and consider $\innerprod[\mathcal{V}/\mathcal{U}]{\mathbf{x}+\mathcal{U},\mathbf{y}+\mathcal{U}}\defeq\innerprod{\mathbf{x},\mathbf{y}}$.
\begin{enumerate}[label=(\roman*)]
\setlength{\itemsep}{0pt}
\item The rule defines a metric form on $\mathcal{V}/\mathcal{U}$ $\Longleftrightarrow$ $\mathcal{U}\leq\operatorname{rad}(\mathcal{V})$.
\item If $\mathcal{U}\leq\operatorname{rad}(\mathcal{V})$, then $\operatorname{rad}(\mathcal{V}/\mathcal{U})=\operatorname{rad}(\mathcal{V})/\mathcal{U}$.
\item The metric space $\mathcal{V}/\operatorname{rad}(\mathcal{V})$ is nonsingular.
\end{enumerate}
\begin{proof}
\begin{enumerate}[label=(\roman*)]
\setlength{\itemsep}{0pt}
\item[$(i)$] ($\Rightarrow$) If the rule defines a metric form, then $(\forall\,\mathbf{u}\in\mathcal{U})(\forall\,\mathbf{y}\in\mathcal{V})\,0=\innerprod[\mathcal{V}/\mathcal{U}]{\mathbf{u}+\mathcal{U},\mathbf{y}+\mathcal{U}}=\innerprod{\mathbf{u},\mathbf{y}}$, so $\mathcal{U}\leq\operatorname{rad}(\mathcal{V})$.
\item[$(i)$] ($\Leftarrow$) Suppose that $\mathcal{U}\leq\operatorname{rad}(\mathcal{V})$. For $\mathbf{u},\mathbf{v}\in\mathcal{U}$ and $\mathbf{x},\mathbf{y}\in\mathcal{V}$, $\innerprod{\mathbf{x}+\mathbf{u},\mathbf{y}+\mathbf{v}}=\innerprod{\mathbf{x},\mathbf{y}}$ by Proposition \ref{prop:metric-radical-equals-bilinear-radicals}(i) and bilinearity.\\
Thus the rule is representative-independent; bilinearity and the form type descend from $\mathcal{V}$.
\item[$(ii)$] Suppose that $\mathcal{U}\leq\operatorname{rad}(\mathcal{V})$. Fix $\mathbf{x}\in\mathcal{V}$. $\mathbf{x}+\mathcal{U}\in\operatorname{rad}(\mathcal{V}/\mathcal{U})\Longleftrightarrow(\forall\,\mathbf{y}\in\mathcal{V})\,\innerprod{\mathbf{x},\mathbf{y}}=0\Longleftrightarrow\mathbf{x}\in\operatorname{rad}(\mathcal{V})$. Hence $\operatorname{rad}(\mathcal{V}/\mathcal{U})=\operatorname{rad}(\mathcal{V})/\mathcal{U}$.
\item[$(iii)$] With $\mathcal{U}=\operatorname{rad}(\mathcal{V})$, item (ii) becomes $\operatorname{rad}(\mathcal{V}/\mathcal{U})=\operatorname{rad}(\mathcal{V})/\operatorname{rad}(\mathcal{V})=\{\mathbf{0}\}$.
\end{enumerate}
\hlt{Structural Commentary}: Item (i) tests representatives; item (ii) identifies the quotient radical; item (iii) removes it.\\
\hlt{Key Insight}: Precisely radical directions may be collapsed without changing pairings.
\end{proof}
\end{proposition}

\begin{proposition} \label{prop:orthogonal-sums-radicals-and-regular-quotients} \hlt{Radicals and Orthogonal Direct Sums}\\
Suppose that $\mathcal{V}=\mathcal{S}\odot\mathcal{T}$, and put $R_{\mathcal{X}}\defeq\operatorname{rad}(\mathcal{X})$ for $\mathcal{X}\in\{\mathcal{V},\mathcal{S},\mathcal{T}\}$.
\begin{enumerate}[label=(\roman*)]
\setlength{\itemsep}{0pt}
\item $R_{\mathcal{V}}=R_{\mathcal{S}}\odot R_{\mathcal{T}}$.
\item $\mathcal{V}/R_{\mathcal{V}}\cong(\mathcal{S}/R_{\mathcal{S}})\odot(\mathcal{T}/R_{\mathcal{T}})$.
\item $\mathcal{V}$ is nonsingular $\Longleftrightarrow$ $\mathcal{S}$ and $\mathcal{T}$ are nonsingular.
\end{enumerate}
\begin{proof}
\begin{enumerate}[label=(\roman*)]
\setlength{\itemsep}{0pt}
\item[$(i)$] For $\mathbf{s}\in\mathcal{S}$ and $\mathbf{t}\in\mathcal{T}$, $\mathbf{s}+\mathbf{t}\in R_{\mathcal{V}}\Longleftrightarrow\mathbf{s}\in R_{\mathcal{S}}\wedge\mathbf{t}\in R_{\mathcal{T}}$ by $\mathcal{S}\perp\mathcal{T}$.\\
Uniqueness and orthogonality of the decomposition establish $R_{\mathcal{V}}=R_{\mathcal{S}}\odot R_{\mathcal{T}}$.
\item[$(ii)$] Define $\Phi((\mathbf{s}+\mathbf{t})+R_{\mathcal{V}})=(\mathbf{s}+R_{\mathcal{S}},\mathbf{t}+R_{\mathcal{T}})$.\\
The map is well-defined by (i), with inverse $\Psi(\mathbf{s}+R_{\mathcal{S}},\mathbf{t}+R_{\mathcal{T}})=(\mathbf{s}+\mathbf{t})+R_{\mathcal{V}}$, and hence is a linear bijection.\\
For $\mathbf{s}_1,\mathbf{s}_2\in\mathcal{S}$ and $\mathbf{t}_1,\mathbf{t}_2\in\mathcal{T}$, $\innerprod{\mathbf{s}_1+\mathbf{t}_1,\mathbf{s}_2+\mathbf{t}_2}=\innerprod{\mathbf{s}_1,\mathbf{s}_2}+\innerprod{\mathbf{t}_1,\mathbf{t}_2}$ by orthogonality. Thus $\Phi$ is an isometry.
\item[$(iii)$] ($\Rightarrow$) If $\mathcal{V}$ is nonsingular, then $R_{\mathcal{S}}\odot R_{\mathcal{T}}=\{\mathbf{0}\}$ by (i), so $\mathcal{S}$ and $\mathcal{T}$ are nonsingular.
\item[$(iii)$] ($\Leftarrow$) If $\mathcal{S}$ and $\mathcal{T}$ are nonsingular, then $R_{\mathcal{V}}=R_{\mathcal{S}}\odot R_{\mathcal{T}}=\{\mathbf{0}\}$ by (i).
\end{enumerate}
\hlt{Structural Commentary}: Item (i) splits the radical; item (ii) descends the splitting; item (iii) tests nonsingularity.\\
\hlt{Key Insight}: Orthogonality separates both pairings and radical directions.
\end{proof}
\end{proposition}

\begin{corollary} \label{cor:metric-spaces-determined-by-regular-quotient-and-dimension} \hlt{Recovering a Singular Space from Its Regular Quotient}\\
If $\dim(\mathcal{V})=\dim(\mathcal{W})$ and $\mathcal{V}/\operatorname{rad}(\mathcal{V})\cong\mathcal{W}/\operatorname{rad}(\mathcal{W})$, then $\mathcal{V}\cong\mathcal{W}$.
\begin{proof}
\begin{enumerate}[label=(\roman*)]
\setlength{\itemsep}{0pt}
\item \hlt{Splitting}: Choose $\mathcal{V}=\operatorname{rad}(\mathcal{V})\odot\mathcal{S}$ and $\mathcal{W}=\operatorname{rad}(\mathcal{W})\odot\mathcal{T}$ with $\mathcal{S},\mathcal{T}$ nonsingular by Theorem \ref{thm:radical-splitting-metric-space}.
\item \hlt{Regular Parts}: The isometries $\mathcal{V}/\operatorname{rad}(\mathcal{V})\cong\mathcal{S}$ and $\mathcal{W}/\operatorname{rad}(\mathcal{W})\cong\mathcal{T}$ hold by Proposition \ref{prop:orthogonal-sums-radicals-and-regular-quotients}(ii).\\
Hence $\mathcal{S}\cong\mathcal{T}$ by the assumed quotient isometry.
\item \hlt{Radicals}: The ambient and regular-part dimension equalities imply $\dim\operatorname{rad}(\mathcal{V})=\dim\operatorname{rad}(\mathcal{W})$.\\
Fix a vector-space isomorphism between the radicals by Theorem \ref{thm:finite-dimensional-isomorphic-iff-same-dimension}; it is an isometry because both restricted forms vanish.
\item \hlt{Assembly}: The orthogonal direct sum of the radical isometry and an isometry $\mathcal{S}\to\mathcal{T}$ is an isometry $\mathcal{V}\to\mathcal{W}$.
\end{enumerate}
\hlt{Structural Commentary}: Radical splittings reduce the classification to regular-part isometry and radical dimension.\\
\hlt{Key Insight}: A zero-form radical is classified solely by dimension.
\end{proof}
\end{corollary}

\begin{definition} \label{def:hyperbolic-pairs-and-spaces} \hlt{Hyperbolic Pairs and Hyperbolic Spaces}
\begin{enumerate}[label=(\roman*)]
\setlength{\itemsep}{0pt}
\item A \hlt{hyperbolic pair} is $(\mathbf{u},\mathbf{v})$ with $\innerprod{\mathbf{u},\mathbf{u}}=\innerprod{\mathbf{v},\mathbf{v}}=0$ and $\innerprod{\mathbf{u},\mathbf{v}}=1$.
\item The plane $\operatorname{span}(\mathbf{u},\mathbf{v})$ of a hyperbolic pair is a \hlt{hyperbolic plane}.
\item A \hlt{hyperbolic space} is an orthogonal direct sum of hyperbolic planes.
\item A basis $(\mathbf{u}_1,\mathbf{v}_1,\ldots,\mathbf{u}_k,\mathbf{v}_k)$ is a \hlt{hyperbolic basis} if each $(\mathbf{u}_i,\mathbf{v}_i)$ is hyperbolic and the planes $\operatorname{span}(\mathbf{u}_i,\mathbf{v}_i)$ are mutually orthogonal.\\
In the symplectic case it is a \hlt{symplectic basis}.
\end{enumerate}
\end{definition}

\begin{theorem} \label{thm:isotropic-vector-sits-in-hyperbolic-plane} \hlt{Capturing Isotropic Vectors}\\
Let $\mathcal{V}$ be nonsingular and either symplectic, or orthogonal with $\operatorname{char}(\mathbb{F})\neq2$.\\
Then $(\forall\,\mathbf{u}\in\mathcal{V}\setminus\{\mathbf{0}\})(\innerprod{\mathbf{u},\mathbf{u}}=0\Longrightarrow(\exists\,\mathbf{v}\in\mathcal{V})\,\innerprod{\mathbf{u},\mathbf{v}}=1\wedge\innerprod{\mathbf{v},\mathbf{v}}=0)$.
\begin{proof}
Fix $\mathbf{u}\neq\mathbf{0}$ with $\innerprod{\mathbf{u},\mathbf{u}}=0$. Since $\mathbf{u}\notin\operatorname{rad}(\mathcal{V})$, fix $\mathbf{x}\in\mathcal{V}$ with $c\defeq\innerprod{\mathbf{u},\mathbf{x}}\neq0$.\\
If $\mathcal{V}$ is symplectic, set $\mathbf{v}=c^{-1}\mathbf{x}$; then $\innerprod{\mathbf{u},\mathbf{v}}=1$ and $\innerprod{\mathbf{v},\mathbf{v}}=0$ by alternation. If $\mathcal{V}$ is orthogonal, set $\mathbf{v}=c^{-1}\mathbf{x}-\frac{\innerprod{\mathbf{x},\mathbf{x}}}{2c^2}\mathbf{u}$.\\
Then $\innerprod{\mathbf{u},\mathbf{v}}=1$ and $\innerprod{\mathbf{v},\mathbf{v}}=c^{-2}\innerprod{\mathbf{x},\mathbf{x}}-c^{-2}\innerprod{\mathbf{x},\mathbf{x}}=0$ by symmetry and $\innerprod{\mathbf{u},\mathbf{u}}=0$. Thus $(\mathbf{u},\mathbf{v})$ is hyperbolic.\\
\hlt{Structural Commentary}: A nonorthogonal partner is rescaled, then corrected by $\mathbf{u}$ only in the orthogonal case.\\
\hlt{Key Insight}: The correction cancels self-pairing without changing the cross-pairing.
\end{proof}
\end{theorem}

\begin{theorem} \label{thm:nonsingular-completions-hyperbolic-extensions} \hlt{Nonsingular Completions}\\
Let $\mathcal{U}\leq\mathcal{V}$, where $\mathcal{V}$ is nonsingular and either symplectic, or orthogonal with $\operatorname{char}(\mathbb{F})\neq2$.\\
Write $\mathcal{U}=\operatorname{span}(\mathbf{u}_1,\ldots,\mathbf{u}_r)\odot\mathcal{W}$, where $(\mathbf{u}_1,\ldots,\mathbf{u}_r)$ is a basis of $\operatorname{rad}(\mathcal{U})$ and $\mathcal{W}$ is nonsingular.
\begin{enumerate}[label=(\roman*)]
\setlength{\itemsep}{0pt}
\item $(\exists\,\mathbf{v}_1,\ldots,\mathbf{v}_r\in\mathcal{V})$ such that $\mathcal{H}_r=\operatorname{span}(\mathbf{u}_1,\mathbf{v}_1,\ldots,\mathbf{u}_r,\mathbf{v}_r)$ is hyperbolic and $\widehat{\mathcal{U}}=\mathcal{H}_r\odot\mathcal{W}$ contains $\mathcal{U}$ with $\dim(\widehat{\mathcal{U}})=\dim(\mathcal{U})+r$.
\item If $\mathcal{U}\leq\mathcal{X}\leq\mathcal{V}$ and $\mathcal{X}$ is nonsingular, then $\dim(\mathcal{X})\geq\dim(\mathcal{U})+r$; equality is attained, and all equality cases are isometric.
\item Let $\widetilde{\mathcal{V}}$ be nonsingular of same kind, $\mathcal{T}\leq\widetilde{\mathcal{V}}$, and $\sigma:\mathcal{U}\xrightarrow{\sim}\mathcal{T}$. Compatible equality-case extensions admit an isometry extending $\sigma$.
\end{enumerate}
\begin{proof}
\begin{enumerate}[label=(\roman*)]
\setlength{\itemsep}{0pt}
\item[$(i)$] The decomposition $\mathcal{V}=\mathcal{W}\odot\mathcal{W}^{\perp}$ holds by Theorem \ref{thm:metric-orthogonal-complement-dimension}(ii), and $(\forall\,i\in\{1,\ldots,r\})\,\mathbf{u}_i\in\mathcal{W}^{\perp}$.\\
Set $\mathcal{H}_0=\{\mathbf{0}\}$ and proceed by induction on $j\in\{1,\ldots,r\}$. Given $\mathcal{H}_{j-1}$, set $\mathcal{S}_j=\operatorname{span}(\mathbf{u}_{j+1},\ldots,\mathbf{u}_r)\odot\mathcal{W}\odot\mathcal{H}_{j-1}$.\\
In a relation $\mathbf{u}_j\in\mathcal{S}_j$, pairing with $\mathbf{u}_1,\ldots,\mathbf{u}_{j-1}$ forces all $\mathbf{v}_i$-coefficients to vanish.\\
Independence of $\mathbf{u}_1,\ldots,\mathbf{u}_r$ modulo $\mathcal{W}$ then contradicts the relation.\\
Hence $\mathbf{u}_j\notin\mathcal{S}_j$. The equality $\mathcal{S}_j^{\perp\perp}=\mathcal{S}_j$ holds by Theorem \ref{thm:metric-orthogonal-complement-dimension}(iii). Since $\mathbf{u}_j\notin\mathcal{S}_j^{\perp\perp}$, fix $\mathbf{x}\in\mathcal{S}_j^{\perp}$ with $c\defeq\innerprod{\mathbf{u}_j,\mathbf{x}}\neq0$.\\
The inclusion $\mathbf{u}_j\in\mathcal{S}_j^{\perp}$ holds by the radical property and the preceding stages.\\
Set $\mathbf{v}_j=c^{-1}\mathbf{x}$ in the symplectic case and $\mathbf{v}_j=c^{-1}\mathbf{x}-\frac{\innerprod{\mathbf{x},\mathbf{x}}}{2c^2}\mathbf{u}_j$ in the orthogonal case.\\
Then $\mathbf{v}_j\in\mathcal{S}_j^{\perp}$, $\innerprod{\mathbf{u}_j,\mathbf{v}_j}=1$, and $\innerprod{\mathbf{v}_j,\mathbf{v}_j}=0$ by alternation or the correction calculation. Put $\mathcal{H}_j=\mathcal{H}_{j-1}\odot\operatorname{span}(\mathbf{u}_j,\mathbf{v}_j)$.\\
After $r$ stages, $\mathcal{H}_r$ is hyperbolic, $\mathcal{U}\leq\widehat{\mathcal{U}}\defeq\mathcal{H}_r\odot\mathcal{W}$, and $\dim(\widehat{\mathcal{U}})=2r+\dim(\mathcal{W})=\dim(\mathcal{U})+r$.
\item[$(ii)$] Choose $\mathcal{X}=\mathcal{U}\oplus\mathcal{Y}$ by Proposition \ref{prop:complements-exist-finite-dimensions}.\\
The pairing $\operatorname{rad}(\mathcal{U})\times\mathcal{Y}\to\mathbb{F}$ has zero left radical: any vector in it pairs trivially with $\mathcal{U}\oplus\mathcal{Y}=\mathcal{X}$ and therefore lies in $\operatorname{rad}(\mathcal{X})=\{\mathbf{0}\}$.\\
Its matrix has rank $r$ by Proposition \ref{prop:bilinear-radicals-matrix-and-dual-maps}(iii), so $\dim(\mathcal{Y})\geq r$ and $\dim(\mathcal{X})=\dim(\mathcal{U})+\dim(\mathcal{Y})\geq\dim(\mathcal{U})+r$.\\
Equality is attained by (i). If equality holds, the construction in (i), performed inside $\mathcal{X}$, produces $\widehat{\mathcal{U}}\leq\mathcal{X}$ with $\dim(\widehat{\mathcal{U}})=\dim(\mathcal{X})$.\\
Hence $\widehat{\mathcal{U}}=\mathcal{X}$ by Proposition \ref{prop:dimension-bounds-and-basis-criteria}(ii). Every equality-case extension is therefore isometric to $\mathcal{W}\odot\mathcal{H}_r$.
\item[$(iii)$] Choose source partners $\mathbf{v}_1,\ldots,\mathbf{v}_r$ by (i). Put $\mathbf{u}'_i=\sigma\mathbf{u}_i$ and $\mathcal{W}'=\sigma\mathcal{W}$, so $\mathcal{T}=\operatorname{span}(\mathbf{u}'_1,\ldots,\mathbf{u}'_r)\odot\mathcal{W}'$.\\
Choose target partners $\mathbf{v}'_1,\ldots,\mathbf{v}'_r$ by (i) in $\widetilde{\mathcal{V}}$, and extend $\sigma$ by $(\forall\,i\in\{1,\ldots,r\})\,\mathbf{v}_i\mapsto\mathbf{v}'_i$.\\
This basis-to-basis map preserves the form on the hyperbolic summands, on $\mathcal{W}\to\mathcal{W}'$, and across the orthogonal summands.\\
It is the required isometry between equality-case extensions.
\end{enumerate}
\hlt{Structural Commentary}: Item (i) pairs radical directions inductively; item (ii) proves minimality; item (iii) transports partners.\\
\hlt{Key Insight}: Each radical basis vector requires exactly one hyperbolic partner.
\end{proof}
\end{theorem}

\begin{definition} \label{def:witt-extension-cancellation-properties} \hlt{Witt Extension and Cancellation Properties}\\
Let $\mathcal{V}\cong\mathcal{W}$ be nonsingular metric vector spaces of the same kind, either symplectic or orthogonal with $\operatorname{char}(\mathbb{F})\neq2$.\\
Write $\tau:\mathcal{A}\xrightarrow{\sim}\mathcal{B}$ when $\tau$ is an isometry. The pair $(\mathcal{V},\mathcal{W})$ has:
\begin{enumerate}[label=(\roman*)]
\setlength{\itemsep}{0pt}
\item the \hlt{Witt extension property} if $(\forall\,\mathcal{S}\leq\mathcal{V})(\forall\,\mathcal{T}\leq\mathcal{W})(\forall\,\tau:\mathcal{S}\xrightarrow{\sim}\mathcal{T})(\exists\,\overline{\tau}:\mathcal{V}\xrightarrow{\sim}\mathcal{W})\,\overline{\tau}|_{\mathcal{S}}=\tau$;
\item the \hlt{Witt cancellation property} if $(\forall\,\mathcal{S}\leq\mathcal{V})(\forall\,\mathcal{T}\leq\mathcal{W})\,[\operatorname{rad}(\mathcal{S})=\operatorname{rad}(\mathcal{T})=\{\mathbf{0}\}\wedge\mathcal{S}\cong\mathcal{T}\Longrightarrow\mathcal{S}^{\perp}\cong\mathcal{T}^{\perp}]$.
\end{enumerate}
\end{definition}

\begin{proposition} \label{prop:witt-extension-cancellation-bridge} \hlt{Witt Extension--Cancellation Bridge}\\
Let $\mathcal{V}\cong\mathcal{W}$ be nonsingular metric vector spaces of the same kind, either symplectic or orthogonal with $\operatorname{char}(\mathbb{F})\neq2$.\\
Then $(\mathcal{V},\mathcal{W})$ has the Witt extension property $\Longleftrightarrow$ it has the Witt cancellation property.
\begin{proof}
\begin{enumerate}[label=\roman*.]
\setlength{\itemsep}{0pt}
\item[$(\Rightarrow)$] Let $\mathcal{S}\leq\mathcal{V}$ and $\mathcal{T}\leq\mathcal{W}$ be nonsingular, fix $\sigma:\mathcal{S}\xrightarrow{\sim}\mathcal{T}$. Decompositions $\mathcal{V}=\mathcal{S}\odot\mathcal{S}^{\perp}$ and $\mathcal{W}=\mathcal{T}\odot\mathcal{T}^{\perp}$ hold by Theorem \ref{thm:metric-orthogonal-complement-dimension}(ii).\\
Fix an extension $\widehat{\sigma}:\mathcal{V}\xrightarrow{\sim}\mathcal{W}$ by Witt extension property. Then $\widehat{\sigma}(\mathcal{S}^{\perp})=\mathcal{T}^{\perp}$ by form preservation and bijectivity, so $\mathcal{S}^{\perp}\cong\mathcal{T}^{\perp}$.
\item[$(\Leftarrow)$] Fix $\sigma:\mathcal{S}\xrightarrow{\sim}\mathcal{T}$ for arbitrary $\mathcal{S}\leq\mathcal{V}$ and $\mathcal{T}\leq\mathcal{W}$.\\
Choose equality-case nonsingular extensions $\widehat{\mathcal{S}}\supseteq\mathcal{S}$ and $\widehat{\mathcal{T}}\supseteq\mathcal{T}$ and an extension $\widehat{\sigma}:\widehat{\mathcal{S}}\xrightarrow{\sim}\widehat{\mathcal{T}}$ by Theorem \ref{thm:nonsingular-completions-hyperbolic-extensions}(iii).\\
The decompositions $\mathcal{V}=\widehat{\mathcal{S}}\odot\widehat{\mathcal{S}}^{\perp}$ and $\mathcal{W}=\widehat{\mathcal{T}}\odot\widehat{\mathcal{T}}^{\perp}$ hold by Theorem \ref{thm:metric-orthogonal-complement-dimension}(ii).\\
Fix $\eta:\widehat{\mathcal{S}}^{\perp}\xrightarrow{\sim}\widehat{\mathcal{T}}^{\perp}$ by the Witt cancellation property. Then $\widehat{\sigma}\odot\eta:\mathcal{V}\xrightarrow{\sim}\mathcal{W}$ extends $\sigma$.
\end{enumerate}
\hlt{Structural Commentary}: Extension identifies complements; cancellation pairs the complements of compatible completions.\\
\hlt{Key Insight}: Nonsingular completion converts arbitrary extension into orthogonal cancellation.
\end{proof}
\end{proposition}

\begin{theorem} \label{thm:witt-extension-theorem-metric-spaces} \hlt{Witt Extension Theorem}\\
Let $\mathcal{V}\cong\mathcal{W}$ be nonsingular metric vector spaces of the same kind, either symplectic or orthogonal with $\operatorname{char}(\mathbb{F})\neq2$.\\
If $\mathcal{S}\leq\mathcal{V}$, $\mathcal{T}\leq\mathcal{W}$, and $\sigma:\mathcal{S}\xrightarrow{\sim}\mathcal{T}$, then $(\exists\,\widehat{\sigma}:\mathcal{V}\xrightarrow{\sim}\mathcal{W})\,\widehat{\sigma}|_{\mathcal{S}}=\sigma$.
\begin{proof}
\begin{enumerate}[label=(\roman*)]
\setlength{\itemsep}{0pt}
\item \hlt{Symplectic Preparation}: Any two nonsingular symplectic spaces of equal dimension are isometric.\\
In dimension zero, the unique linear map is an isometry.\\
At positive dimension, a nonzero vector belongs to a hyperbolic plane by Theorem \ref{thm:isotropic-vector-sits-in-hyperbolic-plane}. The plane splits off by Theorem \ref{thm:metric-orthogonal-complement-dimension}(ii).\\
Its orthogonal complement is nonsingular by Theorem \ref{thm:metric-orthogonal-complement-dimension}(iv).\\
Induction constructs a hyperbolic basis, and matching hyperbolic bases defines the required isometry.
\item \hlt{Orthogonal Line Extension}: Let $\mathcal{V},\mathcal{W}$ be orthogonal, fix $\theta:\mathcal{V}\xrightarrow{\sim}\mathcal{W}$, and suppose $\innerprod{\mathbf{s},\mathbf{s}}=\innerprod{\mathbf{t},\mathbf{t}}\neq0$.\\
Put $\mathbf{x}=\theta\mathbf{s}$, $\mathbf{y}=\mathbf{t}$, and $q(\mathbf{z})=\innerprod{\mathbf{z},\mathbf{z}}$. For $q(\mathbf{a})\neq0$, define $r_{\mathbf{a}}(\mathbf{z})=\mathbf{z}-\frac{2\innerprod{\mathbf{z},\mathbf{a}}}{q(\mathbf{a})}\mathbf{a}$.\\
Direct expansion establishes $\innerprod{r_{\mathbf{a}}\mathbf{z},r_{\mathbf{a}}\mathbf{w}}=\innerprod{\mathbf{z},\mathbf{w}}$, $r_{\mathbf{a}}^2=I_{\mathcal{W}}$, and $r_{\mathbf{a}}\mathbf{a}=-\mathbf{a}$. Thus $r_{\mathbf{a}}$ is isometry. If $q(\mathbf{x}-\mathbf{y})\neq0$, then $r_{\mathbf{x}-\mathbf{y}}\mathbf{x}=\mathbf{y}$.\\
If $q(\mathbf{x}-\mathbf{y})=0$, then $q(\mathbf{x}+\mathbf{y})=4q(\mathbf{y})\neq0$ and $r_{\mathbf{y}}r_{\mathbf{x}+\mathbf{y}}\mathbf{x}=\mathbf{y}$. Thus $r_{\mathbf{x}-\mathbf{y}}\theta$ or $r_{\mathbf{y}}r_{\mathbf{x}+\mathbf{y}}\theta$ extends $\mathbf{s}\mapsto\mathbf{t}$.
\item \hlt{Nonsingular Subspaces}: Suppose that $\mathcal{S},\mathcal{T}$ are nonsingular, and induct on $\dim(\mathcal{S})$.\\
The zero-dimensional case is the assumed isometry $\mathcal{V}\cong\mathcal{W}$.\\
At positive dimension, choose a hyperbolic plane $\mathcal{A}\leq\mathcal{S}$ in the symplectic case by Theorem \ref{thm:isotropic-vector-sits-in-hyperbolic-plane}.\\
In the orthogonal case, fix $\mathbf{s}\in\mathcal{S}$ with $\innerprod{\mathbf{s},\mathbf{s}}\neq0$ by Proposition \ref{prop:form-type-relations-by-characteristic}(ii) and nonsingularity, and set $\mathcal{A}=\operatorname{span}(\mathbf{s})$.\\
In either case, put $\mathcal{B}=\sigma\mathcal{A}$. In the symplectic case, $\mathcal{A}^{\perp}$ and $\mathcal{B}^{\perp}$ are nonsingular and equal-dimensional by Theorem \ref{thm:metric-orthogonal-complement-dimension}(i),(iv).\\
In that case, fix $\beta:\mathcal{A}^{\perp}\xrightarrow{\sim}\mathcal{B}^{\perp}$ by (i), and set $\alpha=(\sigma|_{\mathcal{A}})\odot\beta$. In orthogonal case, fix ambient isometry $\alpha$ extending $\sigma|_{\mathcal{A}}$ by (ii).\\
Set $\mathcal{S}_0=\mathcal{S}\cap\mathcal{A}^{\perp}$ and $\mathcal{T}_0=\mathcal{T}\cap\mathcal{B}^{\perp}$; then $\sigma(\mathcal{S}_0)=\mathcal{T}_0$. Decompositions $\mathcal{S}=\mathcal{A}\odot\mathcal{S}_0$ and $\mathcal{T}=\mathcal{B}\odot\mathcal{T}_0$ hold by Theorem \ref{thm:metric-orthogonal-complement-dimension}(ii).\\
Spaces $\mathcal{S}_0,\mathcal{T}_0$ are nonsingular by Theorem \ref{thm:metric-orthogonal-complement-dimension}(iv). Decompositions $\mathcal{V}=\mathcal{A}\odot\mathcal{A}^{\perp}$ and $\mathcal{W}=\mathcal{B}\odot\mathcal{B}^{\perp}$ hold by Theorem \ref{thm:metric-orthogonal-complement-dimension}(ii).\\
The restriction $\alpha|_{\mathcal{A}^{\perp}}$ identifies the ambient complements, and $\dim(\mathcal{S}_0)<\dim(\mathcal{S})$.\\
Fix $\eta:\mathcal{A}^{\perp}\xrightarrow{\sim}\mathcal{B}^{\perp}$ with $\eta|_{\mathcal{S}_0}=\sigma|_{\mathcal{S}_0}$ by induction. Then $(\sigma|_{\mathcal{A}})\odot\eta:\mathcal{V}\xrightarrow{\sim}\mathcal{W}$ extends $\sigma$.
\item \hlt{General Subspaces}: Choose compatible nonsingular extensions $\widehat{\mathcal{S}},\widehat{\mathcal{T}}$ and $\widehat{\sigma}:\widehat{\mathcal{S}}\xrightarrow{\sim}\widehat{\mathcal{T}}$ extending $\sigma$ by Theorem \ref{thm:nonsingular-completions-hyperbolic-extensions}(iii).\\
An ambient isometry extending $\widehat{\sigma}$ exists by (iii), and it also extends $\sigma$.
\end{enumerate}
\hlt{Structural Commentary}: Reflections or hyperbolic planes drive induction on nonsingular completions of arbitrary subspaces.\\
\hlt{Key Insight}: Compatible completion reduces extension to nonsingular summand induction.
\end{proof}
\end{theorem}

\begin{corollary} \label{cor:witt-cancellation-metric-spaces} \hlt{Witt Cancellation}\\
Let $\mathcal{V}=\mathcal{S}\odot\mathcal{A}$, $\mathcal{W}=\mathcal{T}\odot\mathcal{B}$ be nonsingular metric vector spaces of the same kind, either symplectic or orthogonal with $\operatorname{char}(\mathbb{F})\neq2$.\\
If $\mathcal{V}\cong\mathcal{W}$ and $\mathcal{S}\cong\mathcal{T}$, then $\mathcal{A}\cong\mathcal{B}$.
\begin{proof}
Fix $\sigma:\mathcal{S}\xrightarrow{\sim}\mathcal{T}$, and fix $\widehat{\sigma}:\mathcal{V}\xrightarrow{\sim}\mathcal{W}$ with $\widehat{\sigma}|_{\mathcal{S}}=\sigma$ by Theorem \ref{thm:witt-extension-theorem-metric-spaces}.\\
Identities $\mathcal{A}=\mathcal{S}^{\perp}$ and $\mathcal{B}=\mathcal{T}^{\perp}$ follow from orthogonal decompositions. Hence $\widehat{\sigma}(\mathcal{A})=\widehat{\sigma}(\mathcal{S}^{\perp})=(\widehat{\sigma}\mathcal{S})^{\perp}=\mathcal{T}^{\perp}=\mathcal{B}$, so $\widehat{\sigma}|_{\mathcal{A}}:\mathcal{A}\xrightarrow{\sim}\mathcal{B}$.\\
\hlt{Structural Commentary}: Witt extension sends the first summand to its mate and therefore identifies their complements.\\
\hlt{Key Insight}: An ambient isometry preserves orthogonal complements.
\end{proof}
\end{corollary}

\begin{definition} \label{def:orthogonal-witt-index} \hlt{Witt Index}\\
Let $\mathcal{V}$ be nonsingular orthogonal with $\operatorname{char}(\mathbb{F})\neq2$. Its \hlt{Witt index} is $w(\mathcal{V})\defeq\max\limits\{\dim(\mathcal{L}):\mathcal{L}\leq\mathcal{V},\ \mathcal{L}\subseteq\mathcal{L}^{\perp}\}$.
\end{definition}

\begin{theorem} \label{thm:orthogonal-witt-index-and-anisotropic-decomposition} \hlt{Witt Index and Anisotropic Decomposition}\\
Let $\mathcal{V}$ be orthogonal over $\mathbb{F}$ with $\operatorname{char}(\mathbb{F})\neq2$.
\begin{enumerate}[label=(\roman*)]
\setlength{\itemsep}{0pt}
\item If $\mathcal{V}$ is nonsingular, every maximal totally isotropic subspace has dimension $w(\mathcal{V})$.
\item If $\mathcal{V}$ is nonsingular, every maximal hyperbolic subspace has dimension $2w(\mathcal{V})$.
\item Let $R=\operatorname{rad}(\mathcal{V})$ and $\mathcal{V}=R\oplus\mathcal{U}$ with $\mathcal{U}$ nonsingular.\\
Then $(\exists\,\mathcal{H},\mathcal{S}\leq\mathcal{U})\,\mathcal{H}$ is hyperbolic, $\mathcal{S}$ is anisotropic, $\mathcal{V}=\mathcal{S}\odot\mathcal{H}\odot R$, and $\dim(\mathcal{H})=2w(\mathcal{U})$.
\end{enumerate}
\begin{proof}
\begin{enumerate}[label=(\roman*)]
\setlength{\itemsep}{0pt}
\item[$(i)$] Assume that $\mathcal{V}$ is nonsingular, and let $\mathcal{L}$ be maximal totally isotropic.\\
Choose a totally isotropic $\mathcal{M}$ with $\dim(\mathcal{M})=w(\mathcal{V})$ by the definition of $w(\mathcal{V})$.\\
As $\dim(\mathcal{L})\leq\dim(\mathcal{M})$, choose $\mathcal{M}_0\leq\mathcal{M}$ with $\dim(\mathcal{M}_0)=\dim(\mathcal{L})$. Fix vector-space isomorphism $\tau:\mathcal{L}\to\mathcal{M}_0$ by Theorem \ref{thm:finite-dimensional-isomorphic-iff-same-dimension}.\\
It is an isometry because both restricted forms vanish. Fix $\Phi\in\mathcal{O}(\mathcal{V})$ with $\Phi|_{\mathcal{L}}=\tau$ by Theorem \ref{thm:witt-extension-theorem-metric-spaces}.\\
Then $\mathcal{L}\subseteq\Phi^{-1}\mathcal{M}$, and $\Phi^{-1}\mathcal{M}$ is totally isotropic. Hence $\mathcal{L}=\Phi^{-1}\mathcal{M}$ by maximality, so $\dim(\mathcal{L})=w(\mathcal{V})$.
\item[$(ii)$] Assume that $\mathcal{V}$ is nonsingular. If a hyperbolic subspace has $k$ planes, the span of one first vector from each hyperbolic pair is a $k$-dimensional totally isotropic subspace; hence $k\leq w(\mathcal{V})$. Choose a totally isotropic $\mathcal{M}$ of dimension $w(\mathcal{V})$.\\
Since $\operatorname{rad}(\mathcal{M})=\mathcal{M}$, a hyperbolic completion $\mathcal{H}_0\supseteq\mathcal{M}$ with $\dim(\mathcal{H}_0)=2w(\mathcal{V})$ exists by Theorem \ref{thm:nonsingular-completions-hyperbolic-extensions}(i).\\
Suppose that a maximal hyperbolic $\mathcal{H}$ has dimension $2k<2w(\mathcal{V})$, and choose the span $\mathcal{H}_1\leq\mathcal{H}_0$ of $k$ hyperbolic pairs.\\
Matching hyperbolic bases defines an isometry $\mathcal{H}\to\mathcal{H}_1$. Fix an extension $\Phi\in\mathcal{O}(\mathcal{V})$ by Theorem \ref{thm:witt-extension-theorem-metric-spaces}.\\
Then $\mathcal{H}\subsetneq\Phi^{-1}\mathcal{H}_0$, contrary to maximality. Therefore $\dim(\mathcal{H})=2w(\mathcal{V})$.
\item[$(iii)$] Since $R\perp\mathcal{V}$, the assumed complement satisfies $\mathcal{V}=R\odot\mathcal{U}$.\\
Choose a maximal hyperbolic $\mathcal{H}\leq\mathcal{U}$, and put $\mathcal{S}=\mathcal{H}^{\perp}\cap\mathcal{U}$. The hyperbolic basis makes $\mathcal{H}$ nonsingular.\\
The identities $\dim(\mathcal{H})=2w(\mathcal{U})$ and $\mathcal{U}=\mathcal{H}\odot\mathcal{S}$ hold by (ii) and Theorem \ref{thm:metric-orthogonal-complement-dimension}(ii). Space $\mathcal{S}$ is nonsingular by Theorem \ref{thm:metric-orthogonal-complement-dimension}(iv).\\
If $\mathcal{S}$ contained a nonzero isotropic vector, a hyperbolic plane $\mathcal{P}\leq\mathcal{S}$ would contain it by Theorem \ref{thm:isotropic-vector-sits-in-hyperbolic-plane}.\\
Then $\mathcal{H}\odot\mathcal{P}$ would contradict maximality of $\mathcal{H}$. Thus $\mathcal{S}$ is anisotropic and $\mathcal{V}=\mathcal{S}\odot\mathcal{H}\odot R$.
\end{enumerate}
\hlt{Structural Commentary}: Witt extension fixes maximal isotropic and hyperbolic dimensions; splitting isolates the anisotropic residue.\\
\hlt{Key Insight}: Each hyperbolic plane contributes one maximal isotropic dimension.
\end{proof}
\end{theorem}
